%!TEX root = ./poly.tex
%
\section{Introduction}
\label{sec:intro}

This paper establishes some foundational computations for studying $p$-adic cohomology through the lens of genuine equivariant homotopy theory. Recall that the interface between these is topological Hochschild homology (here viewed as a proper-genuine $\T$-spectrum), and its various types of fixed point spectra. For nice input rings, the homotopy groups of these are given by important arithmetic invariants, often with additional \emph{filtrations} and \emph{twists}.

\begin{example}
If $R$ is a quasiregular semiperfectoid ring, the ($p$-completed) $\TCmin$ and $\TP$ of $R$ are related to its (Nygaard-filtered, Nygaard-completed, Breuil-Kisin twisted, absolute) prismatic cohomology via
\begin{align*}
  \TCmin_{2i}(R;\Z_p) &= \Nyg^{\ge i}\hat\prism_R\{i\}\\
  \TP_{2i}(R;\Z_p) &= \hat\prism_R\{i\}
\end{align*}
(and the odd groups vanish) by work of Bhatt-Morrow-Scholze \cite{BMS2}. In the $q$-de Rham case, the twist can be understood as tensoring with a free module on the symbol $(p^\infty)_q = \frac{q^{p^\infty}-1}{q-1} = \frac{\log_\prism(q^p)}{q-1}$ (see Notation \ref{not:canonical} and \cite[Proposition 2.6.1]{APC}). Antieau-Krause-Nikolaus have announced a generalization of this result in terms of their prismatic cohomology relative to a $\delta$-ring $A$ \cite{RelativePrismatic} whenever $A$ descends along $\S_p\to\Z_p$.
\end{example}

\begin{example}
  When $S$ is $p$-completely smooth over any perfectoid ring $R$ (i.e.\ the $p$-completion of a smooth $R$-algebra), Hesselholt \cite{LarsPTyp} showed that $\TR^n(S;\Z_p)$ is given by the $p$-complete relative de Rham-Witt complex \cite{LangerZinkDRW} adjoined with a polynomial generator in degree 2:
  \begin{equation}
    \label{eq:hesselholt}
    \TR^n_*(S;\Z_p) \isom \dRW[n]^*_{\Spf S/\Spf R}[\sigma_n]\qquad |\sigma_n|=2
  \end{equation}
  Here the $p$-complete de Rham-Witt complex is
  \[ \dRW[n]^*_{\Spf S/\Spf R} \defeq \lim_s \dRW[n]^*_{(S/p^s)/(R/p^s)} \]
  as in \cite[Definition 10.11]{BMS1}; we will subsequently abuse notation and simply write this as $\dRW[n]^*_{S/R}$.
  The identification \eqref{eq:hesselholt} can be made canonical by interpreting $\sigma_n$ as tensoring with the symbol $\xi_{n+1}$ ($= (p^{n+1})_{q^{1/p^{n+1}}}$ in the $q$-de Rham case), see Notation \ref{not:canonical}.

  In particular, the Restriction maps on the twists get multiplied by powers of $\xi_{n+1}/\xi_n=\phi^{-n}(\xi)$, so there are potentially $\plainlim^1$ terms when we take the limit over this result to understand $\TR(S;\Z_p)$. Fortunately, these vanish in the main cases of interest: if $R=k$ is a perfect $\F_p$-algebra, then
  \[ \TR_*(S) = \dRW^*_S, \]
  and if $R=\O_C$ is the ring of integers of a spherically complete perfectoid field $C$, then we have
  \[ \TR_*(S;\Z_p) = \dRW^*_{S/\O_C}[\beta], \qquad |\beta| = 2 \]
   by \cite[Lemma 7.22]{MathewK1}. This identification can be made canonical by interpreting $\beta$ as tensoring with the ideal $(q-1)$ of $\Ainf(\O_C)$.
\end{example}

By Kan extending the Postnikov filtration (or its double-speed acceleration), one gets filtrations on these spectra for arbitrary animated rings, whose associated gradeds are given by derived versions of the same arithmetic invariants.

The genuine equivariant\footnote{$\RO(G)$-graded homotopy exists already for Borel $G$-spectra, but the notion is emphasized less there.} perspective provides a significant enlargement of the homotopy groups of a spectrum: a $G$-spectrum $E$ naturally has homotopy groups $\pi_V^H E$ for each subgroup $H\le G$ and each real virtual representation $V\in\RO(H)$ of $H$. This has a great deal of functoriality in $H$, leading to \emph{$\RO(G)$-graded homotopy Mackey functors $\mpi_\rog E$}. When $E$ is a $G$-$\E_\infty$-ring, we even get an \emph{$\RO(G)$-graded Tambara functor} in the sense of Angeltveit-Bohmann \cite{GradedTambara}.

\begin{remark}
  For all notation and background, the reader can consult
  \begin{itemize}
    \item \cite[Part 1]{SulSliceTHH} and \cite{HillHandbook} for equivariant homotopy theory in general;
    \item \cite[\sec4.2]{SulSlopes} for poly-representations;
    \item \cite[\sec2]{PrismsTambara1} for Tambara functors;
    \item \sec\ref{sec:notation} for all the notation used in this paper.
  \end{itemize}
  In particular, Notation \ref{not:canonical} elaborates what we mean by ``canonical''. In short, $\RO(G)$-graded homotopy groups naturally have the structure of \emph{spans} relative to the $\Z$-graded homotopy groups, and it is as such that we identify them.
\end{remark}

In this paper, we will generalize the above two computations above to $\RO(G)$-grading. Here are some reasons for doing so.

\begin{motivation}
  In general, the $\RO(G)$-graded homotopy groups of a $G$-spectrum tend to give a simpler or more holistic picture, in which the $\Z$-graded groups sit as a funny diagonal. For example, while the integer-graded $\TCmin$ mixes the Nygaard filtration with Breuil-Kisin twists, the $\RO(\T)$-graded groups will give us the Nygaard filtration on the nose:
  \[
    \TCmin_{2i}(R;\Z_p) = \Nyg^{\ge i}\hat\prism_R\{i\}
    \quad\text{but}\quad
    \TCmin_{i\lambda}(R;\Z_p) = \Nyg^{\ge i}\hat\prism_R.
  \]

  This example showcases another important point: the $\RO(G)$-graded groups considered here do not yield ``new'' arithmetic invariants. Rather, they vastly increase the supply of additional \emph{filtrations} and \emph{twists} that can be passed through the topological dictionary.
\end{motivation}

\begin{motivation}
  The original motivation for studying $\RO(G)$-graded $\TR$ specifically comes from algebraic $K$-theory: for any commutative ring $A$, we may consider the reduced $p$-adic $K$-theory of truncated polynomial algebras over $A$,
  \[ \~\K(A[x]/x^e;\Z_p) \defeq \fib\Big( \K(A[x];\Z_p) \to \K(A;\Z_p) \Big) \]
  By \cite[Proposition 8]{HesselholtHandbook}, this can be understood in terms of ``$\RO(G)$-graded big $\TR$'': there is a fiber sequence
  \begin{equation}
    \label{eq:syntomic-drw}
    \~\K(A[x]/x^e;\Z_p) \to \bigTR(\Sigma^{\hsl 1}\THH(A;\Z_p)) \to \bigTR(\Sigma^{\hsl{1/e}}\THH(A;\Z_p))
  \end{equation}
  and the latter two terms are extremely computable. The notation $\hsl\nu$ is recalled in Definition \ref{def:regular-poly}.

  A basic computation is that for $R$ perfectoid,
  \[ \bigTR_{2i-\hsl\nu}(R; \Z_p) = \W_{\lf i/\nu\rf}(R); \]
  this is essentially \cite[Proposition 9.1]{HMFinite} + \cite[Theorem 4.2.10]{HMCyclicPolytopes}, but see \cite[Corollary 4.6]{SulSlopes} for the present formulation (which is mildly more general). This immediately implies the syntomic cohomology
  \[ \~{\Z_p(i)}(R[x]/x^e) = \frac{\W_{ei}(R)}{\V^e\W_i(R)}[-1]. \]

  In \cite{SulSlopes}, this syntomic cohomology is computed by first computing $\Nyg^{\ge *}\prism_{k[x]/x^e}$ and $\prism_{k[x]/x^e}$, but the big Witt vectors appear somewhat miraculously. This method makes the appearance of big Witt vectors clearer, and also generalizes to perfectoid rings. The calculation of $\~\Z_p(i)(R[x]/x^e)$ for $R$ perfectoid was done independently by Riggenbach  using different methods \cite[Theorem 1.1]{RiggenbachTruncated}, and we will build on some of his tools in \sec\ref{sec:qrsp}.

  We therefore view this method of computing syntomic cohomology of monoid algebras as superior to the one used in \cite[\sec3]{SulSlopes}. The main disadvantage is that it is extremely difficult to establish the fiber sequence \eqref{eq:syntomic-drw} \cite{HMCyclicPolytopes}. However, this should be much easier to do in the modern prismatic approach, and to extend to more general monoid algebras. We will not pursue this here, but our work is intended to be a repository of computations to accompany such an endeavor.

  The $\RO(G)$-graded $\TR$ groups relevant to algebraic $K$-theory are the ``negative range'' (integer minus actual representation), which is the easiest range to compute. In this range, the answer has in some sense been known since the author was two and a half years old:
  \begin{itemize}
    \item the base case was proved in \cite[Proposition 9.1]{HMFinite}, and re-proven in \cite[Theorem 8.3]{ROS1TR} from the perspective of the homotopy orbits to TR spectral sequence (our main tool in this paper);
    \item the base case essentially implies the result for smooth algebras over it by \cite[Theorem 11]{HesselholtHandbook};
    \item the perfectoid base case was done in \cite{SulSlopes}---but in this range it behaves exactly like the characteristic $p$ case.
  \end{itemize}
  However, it took the author over thirty years to understand this.

  Even in the negative range, we get two new things. Our canonical identifications make it possible to compute the functoriality of $\~\K(A[x]/x^e;\Z_p)$ (or more general monoid algebras) in $e$. We also obtain an interesting filtration on these groups by computing with cell structures instead of the homotopy orbits to $\TR$ spectral sequence. These applications will be deferred to a companion paper.
\end{motivation}

While the $\RO(G)$-graded groups themselves are not ``new'', they provide a home for something which is.

\begin{motivation}[Tambara structure]
  Given $K\le H\le G$, there are $\RO(G)$-graded Frobenius (restriction) and Verschiebung (transfer) maps that go
    \[\xymatrix{
        \pi^H_V E \ar@<0.5ex>[r]^-F & \pi^K_{\res^H_K(V)} \ar@<0.5ex>[l]^-V E
    }\]
    Since the restriction of a trivial representation is trivial, these restrict to $F$ and $V$ maps on the $\Z$-graded homotopy groups.

    Genuine equivariant homotopy theory emphasizes another map, the \emph{Norm\iffalse\footnote{Frobenius is capitalized after German mathematician Ferdinand Frobenius, Verschiebung is capitalized because it's a German noun, and people capitalize Restriction out of consistency. The Norm map is named for Canadian mathematician Norm Macdonald.}\fi{} map}, which goes
    \[ N^H_K\colon \pi^K_V E \morph \pi^H_{\ind^H_K(V)} E \]
    when $E$ is a $G$-$\E_\infty$-ring spectrum. Since the induction of a trivial representation is never trivial, except for the zero representation, studying the norm in nonzero $\Z$-graded degrees forces us to contemplate $\RO(G)$-graded degrees. Understanding these norm maps was our primary motivation for these calculations.

  As for the zero representation, the $\pi_0$ Norm was introduced in \cite{AngeltveitNorm} for Witt vectors (and therefore perfect prisms), and for transversal prisms in \cite{PrismsTambara1}. Two applications of the latter are a conceptual construction of de Rham-Witt comparison maps for prismatic cohomology \cite{Molokov}, \cite[\sec4.1]{PrismsTambara1}, and (at odd primes) a refinement of the prismatic logarithm \cite{MaoLogarithm}. It would be very interesting to try to generalize these applications to $\RO(G)$-graded norms.
\end{motivation}

\begin{warning}
  Unfortunately, much of the literature refers to the sum-over-conjugates map
  \[ X_{hG} \to X^{hG} \]
  as the norm, where $X_{hG}$ denotes the homotopy orbits in spectra. By analogy with Galois theory, this should really be called the \emph{trace}. The true norm is the analogue of the above where the orbits are taken in $\E_\infty$-monoids and we take a product over conjugates rather than a sum.
\end{warning}

\begin{motivation}
  $\RO(G)$-graded homotopy groups also serve as calculational input for studying the equivariant slice filtration. Understanding the role of norms is also very important for understanding the slice filtration, since it is tailor-made to interact well with them. The author's original entrypoint into this subject was to attempt to modify the BMS filtration by replacing the Postnikov filtration with the slice filtration, which amounts to studying the slice filtration on $\TCmin(S;\Z_p)$ and $\TP(S;\Z_p)$ when $S$ is quasiregular semiperfectoid. This was our primary motivation for the calculations in \sec\ref{sec:qrsp}, and is now within reach.
\end{motivation}

\subsection{Borel spectra and prismatic cohomology}
We first compute the $\RO(\T)$-graded $\TCmin$ and $\TP$ of quasiregular semiperfectoid rings (\sec\ref{sec:qrsp}). Let us quickly review the representation theory of $\T$ and set some notation. The complex representation ring of $\T$ is $\RU(\T)=\Z[\lambda^\pm]$, with $\lambda$ the natural representation of $\T$ on $\C$; explicitly, $\lambda^i=\C$ with $z\in\T$ acting as multiplication by $z^i$. Complex conjugation induces a real isomorphism $\lambda^i \isom \lambda^{-i}$, so the real representation ring is $\RO(\T)=\Z[\lambda]$. Since we are only concerned with the homotopy type of representation spheres, we will work with $\Z[\lambda]$; but $\lambda^0$ will still mean $\C^{\mathrm{triv}}$ rather than $\R^{\mathrm{triv}}$. In the $p$-typical context, we have $S^{\lambda^i} \simeq S^{\lambda^j}$ iff $v_p(i)=v_p(j)$, so we set $\lambda_i=\lambda^{p^i}$ and $\lambda_\infty=\lambda^0$ (so we should really say $\RO(C_{p^\infty})$-graded). 
In addition to the irreducible decomposition
\[ \alpha = k_0 \lambda_0 + k_1 \lambda_1 + \dotsb + k_n \lambda_n + k_\infty \lambda_\infty, \]
it will be convenient to encode representations via the eventually-constant sequence
\[ d_n(\alpha) = \dim_\C \alpha^{C_{p^n}} \]
which we write as $(d_0,\dotsc,d_n;d_\infty)$. This encoding was introduced by Gerhardt and Angeltveit-Gerhardt \cite{TeenaTR,ROS1TR}, and as we will see, it is suited to working in ``ghost coordinates''.

Recall the ideals
\[ I_n \defeq I\phi(I)\dotsm\phi^{n-1}(I) \]
from \cite[Notation 2.2.2]{APC} (warning: our notation is off-by-one from that used in \cite{PrismsTambara1}).

\begin{theorem}
  \label{thm:qrsp}
  Let $A$ be a $\delta$-ring admitting a spherical lift $\S_A$, and let $S$ be a relatively quasiregular semiperfectoid $A$-algebra with associated prism $(\prism_{S/A},I)$. Let
  \begin{align*}
    \alpha
    &=k_0\lambda_0 + \dotsb + k_n\lambda_n + k_\infty \lambda_\infty\\
    &= (d_0,\dotsc,d_n;d_\infty)
  \end{align*}
  be a virtual $C_{p^\infty}$-representation. Then there are canonical identifications
  \begin{align*}
    \TCmin_\alpha(S\rel\S_A)
      &= \Nyg^{\ge d_0}\hat\prism^{(1)}_{S/A}\{k_\infty\} \tnsr \bigotimes_{i=1}^n I_i^{k_i}
    & \TP_\alpha(S\rel\S_A)
    &= \hat\prism^{(1)}_{S/A}\{k_\infty\} \tnsr \bigotimes_{i=1}^n I_i^{k_i}\\
    &= \hat\prism^{(1)}_{S/A} \tnsr \bigotimes_{i=0}^\infty \phi^{i-1}(I^{d_i})
    &&= \hat\prism^{(1)}_{S/A} \tnsr \bigotimes_{i=1}^\infty \phi^{i-1}(I)^{d_i}
  \end{align*}
  To interpret the second line: we have absorbed the Breuil-Kisin twist into the product via $\prism\{1\} = \bigotimes_{i=1}^\infty \phi^{i-1}(I)$, and note that $\phi^{i-1}(I^{d_i}) = \phi^{i-1}(I)^{d_i}$ for $i>0$.
\end{theorem}

\begin{warning}
  Given a $G$-spectrum $E$, there are two notations for the integer-graded homotopy groups of $E^G$: either $\pi_n E^G$ or $E^G_n$. We can imagine these as
  \begin{align*}
    \pi_n E^G &\defeq [S^n, E^G]_\Sp\\
    E^G_n &\defeq [S^n, E]_{\Sp^G}
  \end{align*}
  In other words, for the first version we take $G$-fixed points of $E$ and then take $\pi_n$ of the resulting ordinary spectrum, while in the second we compute the mapping space in $\Sp^G$ directly.

  In the $\RO(G)$-graded context, the first version \emph{does not make sense}: for $\alpha\in\RO(G)$, we define
  \[ E^G_\alpha \defeq [S^\alpha, E]_{\Sp^G} \]
  (where $S^\alpha$ is the one-point compactification of $\alpha$ when $\alpha$ is an actual representation), and this cannot be recovered from the plain spectrum $E^G$. We will often abuse notation by writing $\pi_\alpha E^G$, but it should be remembered that this is technically incorrect.

  Thus, when we write $\TCmin_\alpha$ and $\TP_\alpha$, these are defined as
  \begin{align*}
    \TCmin_\alpha &\defeq [S^\alpha, \THH^{E\T_+}]_{\Sp^\xi} &
    \TP_\alpha &\defeq [S^\alpha, \widetilde{E\T}\tnsr \THH^{E\T_+}]_{\Sp^\xi}\\
    &= \pi_0 (\Sigma^{-\alpha} \THH)^{h\T} &
    &= \pi_0 (\Sigma^{-\alpha} \THH)^{t\T}
  \end{align*}
  with $\Sp^\xi$ the category of cyclonic\footnote{We use the notation $\Sp^\xi$ because $\xi$ looks like a tornado.} spectra (i.e.\ proper-genuine $\T$-spectra).
\end{warning}

% We can Kan extend this computation to get a statement for more general rings. Repeating the sheafiness arguments of \cite[\sec4]{RiggenbachTruncated}, our calculations imply
%
% \begin{corollary}
%   \label{cor:qsyn}
%   Let 
%   \begin{align*}
%     \alpha
%     % &=k_0\lambda_0 + \dotsb + k_n\lambda_n + k_\infty \lambda_\infty\\
%     &= (d_0,\dotsc,d_n;d_\infty)
%   \end{align*}
%   be a virtual $\T$-representation. Then there are equivalences of quasisyntomic sheaves
%   \begin{align*}
%     \gr^i(\Sigma^{-\alpha} \THH(-;\Z_p))^{h\T}
%       &= \left(\bigotimes_{r=1}^n (\varphi^{r-1})^*(\I)^{\tnsr(d_r-d_\infty)}\right)\tnsr\Nyg^{\ge d_0 + i}\O_{\hat\prism}\{d_\infty + i\}[2i]\\
%     \gr^i(\Sigma^{-\alpha} \THH(-;\Z_p))^{t\T}
%       &= \fade{\left(\bigotimes_{r=1}^n (\varphi^{r-1})^*(\I)^{\tnsr(d_r-d_\infty)}\right)\tnsr{}}\O_{\hat\prism}\fade{\{d_\infty+i\}[2i]}\\
%       \gr^i(\Sigma^{-\alpha} \THH(-;\Z_p))^{hC_{p^n}}
%       &= \fade{\left(\bigotimes_{r=1}^n (\varphi^{r-1})^*(\I)^{\tnsr(d_r-d_\infty)}\right)\tnsr{}}\frac{\Nyg^{\ge d_0 + i}\O_{\hat\prism}}{\Nyg^{\ge d_0 + i + 1}\O_{\hat\prism}\tnsr \I_n}\fade{\{d_\infty + i\}[2i]}\\
%       \gr^i(\Sigma^{-\alpha} \THH(-;\Z_p))^{tC_{p^n}}
%       &= \fade{\left(\bigotimes_{r=1}^n (\varphi^{r-1})^*(\I)^{\tnsr(d_r-d_\infty)}\right)\tnsr{}}\frac{\O_{\hat\prism}}{\I_n}\fade{\{d_\infty+i\}[2i]}
%   \end{align*}
% \end{corollary}

The proof of Theorem \ref{thm:qrsp} is very easy, and the argument works much more generally; all we need to do is formulate it. Recall the following generalization of cyclotomic polynomials.
\begin{definition}[{\cite[Remark 2.4.3]{GNS}}]
  \label{def:fgl-cyclotomic}
  Let $\G$ be a formal group law with $n$-series $[n]_\G(t)$. Define
  \[ \Phi_{n,\G}(t) = \prod_{d\mid n}[d]_\G(t)^{\mu(n/d)} \]
  where $\mu$ is the M\"obius function. Note that by M\"obius inversion we have
  \begin{align*}
    [n]_\G(t) &= \prod_{d\mid n}\Phi_{d,\G}(t)
  \end{align*}
\end{definition}

\begin{theorem}[Remark \ref{rmk:even}]
  \label{thm:even}
  Let $E$ be a $\T$-equivariant even $\E_\infty$-ring with associated formal group $\G$. Let $A\{i\} = \pi_{2i} E^{t\T}$, and let $\Fil^\bullet A\{i\}$ be the abutment filtration of the Tate spectral sequence. Fix a choice of complex orientation $t\in A\{-1\}$, define $\<n\>_\G=\frac{[n]_\G(t)}t\in A$ and $\Phi_{n,\G}=\Phi_{n,\G}(t)\in A\ (n>1)$ as above, and assume that $\<n\>_\G$ is a non-zerodivisor in $A$ for all $n$.

  Let $\alpha$ be a virtual complex $\T$-representation with irreducible decomposition
  \[
    \alpha = \sum_{i=0}^\infty k_i \lambda^i
  \]
  and dimension sequence $d_n = \dim_\C \alpha^{C_n},\, d_\infty = \dim_\C \alpha^\T$.

  Then the $\RO(\T)$-graded homotopy groups of $E^{h\T}$ and $E^{t\T}$ are given canonically by
  \begin{align*}
    E^{h\T}_\alpha
      &= \Fil^{d_1}A\{k_0\}\tnsr_A \bigotimes_{n>1} (\<n\>_\G A)^{\tnsr k_n}\\
      &= \Fil^{d_1} A \tnsr_A \bigotimes_{n>1} (\Phi_{n,\G} A)^{\tnsr d_n}\\
      &\fade{\defeq \Fil^{d_1} A\{d_\infty\} \tnsr_A \bigotimes_{n>1} (\Phi_{n,\G} A)^{\tnsr (d_n-d_\infty)}}\\
    E^{t\T}_\alpha
      &= A\{k_0\}\tnsr_A \bigotimes_{n>1} (\<n\>_\G A)^{\tnsr k_n}\\
      &= \bigotimes_{n>1} (\Phi_{n,\G} A)^{\tnsr d_n}\\
      &\fade{\defeq A\{d_\infty\} \tnsr_A \bigotimes_{n>1} (\Phi_{n,\G} A)^{\tnsr (d_n-d_\infty)}}
  \end{align*}
\end{theorem}

We can get more interesting invariants by going ``more genuine'' and considering the ``higher topological negative cyclic homology'' spectra $\TC^{(-n)} \defeq (\TR^n)^{h\T/C_{p^n}}$. We study this in \sec\ref{sub:iterated-nygaard}.


\subsection{Polygonic spectra and de Rham-Witt complexes}
We next shift our attention to $\TR$. Since the Restriction maps involve geometric fixed points, $\TR$ is not graded by representations of $C_{p^\infty}$ or $\T$, but rather \emph{poly-representations} of $\Z_p$ or $\widehat\Z$ (introduced in \cite[\sec4.2]{SulSlopes}). A poly-representation $\m\alpha$ consists of representations $\m\alpha(\Z/n)\in\RO(\Z/n)$ for each $n$, compatible under the maps
\begin{align*}
  \RO(\Z/mn) &\to \RO\left(\frac{\Z/mn\Z}{m\Z/mn\Z}\right)\\
              &= \RO(\Z/m)\\
  \alpha &\mapsto \alpha^{m\Z/mn\Z}
\end{align*}
The name \emph{poly-}representation is chosen partly because it is a family of many representations, but more importantly to emphasize the connection with polygonic spectra \cite{KMNPolygonic}: the representation spheres $S^{\m\alpha(\Z/n)}$ assemble into a polygonic spectrum $S^{\m\alpha}$ (which in particular is $\tnsr$-invertible). We will mostly work with $p$-typical poly-representations, in which case $\widehat\Z$ is replaced by $\Z_p$. We will study the polyrepresentation-graded $\TR$ of smooth algebras $S$ over a perfectoid base $R$.

\begin{remark}
Say that a continuous representation $V$ of $\widehat\Z$ is \emph{quasifinite-dimensional}\footnote{Of course it doesn't make sense for a dimension to be quasifinite; this term is chosen for consistency with the \emph{quasifinitely genuine} $\Z$-spectra from \cite{KMNPolygonic}.} if
\[ V(\Z/n) \defeq V^{n\widehat\Z} \fade{{}\curvearrowleft \Z/n} \]
is finite-dimensional for all $n>0$ (but $V^e$ could be countably-infinite-dimensional). Then a poly-representation is equivalently an element of the Grothendieck group $\RO(\widehat\Z)^\land$ of quasifinite-dimensional $\widehat\Z$-representations, a completion of the usual representation ring $\RO(\widehat\Z)$.
\end{remark}

Flipping and reversing the dimensional encoding, we write a $p$-typical poly-representation as
\[ \m \alpha = (\m d_0, \m d_1, \dotsc), \]
where $\m d_i=\dim_\C\m\alpha(C_{p^n})$ should be thought of $\dim_\C(\text{``$\alpha^{1/C_{p^i}}$''})$, and indeed we sometimes write $d_{-i}(\m\alpha)$ instead of $\m d_i$.

% \begin{remark}
%   The truncated polynomial algebra $S[x]/x^e$ can also be written as $S[y^{1/e}]/y$, where $y=x^e$. This perspective is fruitful as the the limit $R[y^{1/p^\infty}]^\land_p$ is an important example of a quasi-regular semiperfectoid ring. The poly-representation perspective suggests that we should expand ``truncated polynomial algebras'' to include $S[y^\nu]/y$ for general $\nu\in\Q_{>0}$. We expect (but have not checked) that our results hold for such rational-slope truncated polynomial algebras.
% \end{remark}

The pioneering work on $\RO(G)$-graded $\TR$ is by Gerhardt and Angeltveit-Gerhardt \cite{TeenaTR,ROS1TR}, who gave an algorithm to compute $\TR^n_\rog(k)$ as $W_n(k)$-modules when $k$ is a perfect $\F_p$-algebra. We revisited their work in \cite{SulROGTHH} to compute the full rings $\TF_\rog(k)$ and $\TF_\rog(R^+;\Z_p)$ when $k$ is a perfect $\F_p$-algebra and $R^+$ is a $p$-torsionfree perfectoid ring. This in principle implies the computation of $\TR_\rog$ via the long exact sequences
\[ \TF_{\alpha+\lambda_i}(R;\Z_p) \xra{a_{\lambda_i}} \TF_\alpha(R;\Z_p) \to \dbox{$\TR^i_\alpha(R;\Z_p)$} \xra\partial \TF_{\alpha+\lambda_i-1}(R;\Z_p), \]
but it is nontrivial to unpack this; we will instead proceed directly by adapting the homotopy orbits to TR spectral sequence. The results of \cite{SulROGTHH} may be informally summarized as follows:
\begin{itemize}
  \item in even degrees, $\TF_\rog(R;\Z_p)$ is torsionfree and satisfies an $\RU(\T)$-graded form of B\"okstedt periodicity;
  \item in the characteristic $p$ case, there are additional classes witnessing the ``collision'' of the distinguished elements $d=\phi(d)=\phi^2(d)=\dotsb$;
  \item the odd degrees contain deep caverns with dangerous monsters\footnote{Infamously, \cite{SulROGTHH} contains an error related to extension problems in these odd degrees. The author now knows how to fix this and will update the paper soon. This error does not affect any of the computations in this paper on account of restricting to the contiguous range.}.
\end{itemize}

In light of the extraordinary complexity of even the base case, we will not attempt to compute the full $\TR_\rog(S)$, but rather focus our attention on certain safe regions:
\begin{definition}
  Consider the following swathes of poly-representations.

  \begin{itemize}
    \item the \emph{negative range} consists of $\rog=*-V$ with $*\in\Z$ and $V$ an actual poly-representation; equivalently, $\m\alpha=(\m d_0,\m d_1,\dotsc)$ is in the negative range if $\m d_i \ge \m d_{(i+1)}$ for all $i$. This is the range that arises in algebraic $K$-theory calculations. 
    \item the \emph{positive range} consists of $\rog=V-*$ with $*\in\Z$ and $V$ an actual poly-representation; equivalently, $\m\alpha=(\m d_0,\m d_1,\dotsc)$ is in the positive range if $\m d_i \le \m d_{(i+1)}$ for all $i$. This range is of interest because the Norm from the $\Z$-graded de Rham-Witt complex lands in here, allowing us to extend the Witt vector Norm \cite{AngeltveitNorm,PrismsTambara1} to the de Rham-Witt complex. (Also, a tiny portion of it is used in computing the slice filtration \cite[Lemma 4.11]{SulSliceTHH}.)
  \item the \emph{contiguous range} consisting of those $\m\alpha=(\m d_0,\m d_1,\dotsc)$ for which $\{i \mid d_i(\m\alpha)\ge 0\}$ is a contiguous interval $[-\ell, -k]$ or $(-\infty, k]$. We say that $\m\alpha$ is \emph{non-negatively concentrated in $[-\ell, -k]$} or $(-\infty, k]$.
  \end{itemize}
\end{definition}

The negative and positive ranges are exactly the ranges that can be handled by the cellular method; however, this is a very difficult way to proceed. We instead skip straight to the contiguous range and employ the \emph{homotopy orbits to TR spectral sequence} (HOTRSS) introduced in \cite{ROS1TR} and upgraded in \cite{SulROGTHH}. The homotopy orbits to TR spectral sequence can handle all of $\TR_\rog$, but the answer becomes unwieldy outside of the contiguous range.

We need two more ingredients to state the result. Recall that the \emph{canonical filtration} of the de Rham-Witt complex is given by
\begin{align*}
  \Fil^{k+1} \dRW[n]^j_{S/R}
    &= \ker(\dRW[n]^j_{S/R} \xra{R^{n-k}} \dRW[k]^j_{S/R})\\
    &= \im(V^{k+1} + dV^{k+1})
\end{align*}
for $n\in\N\cup\{\infty\}$. (Remember that in our notation Witt vectors start at $\rmW_0$.) Next, in the $\Z$-graded case, Hesselholt's theorem \cite{LarsPTyp} says that $\TR_*$ splits into contributions from individual $\dRW[\bullet]^j$ terms. This is also true in the fancy-graded case (either from the proof of Hesselholt's theorem, or by the homotopy orbits to TR spectral sequence), so we let $\TR_{\rog,j}(S)$ denote the $\dRW^j$ contribution to $\TR_{\rog+j}(S)$. Pulling out $j$ this way, the computation becomes essentially the same for the base case $S=R$, although we have to be slightly more careful with $\Fil$.

\begin{theorem}[Contiguous range of $\TR_\rog$, Theorem {\ref{thm:tr-ctg}}]
  \label{thm:intro-tr-ctg}
  Let $\m\alpha = (\m d_0,\dotsc, \m d_n)$ be a $\Z/p^n$-polyrepresentation which is non-negatively concentrated in $[-\ell, -k]$.

  If $S$ is $p$-completely smooth over a $p$-torsionfree perfectoid ring $R^+$, then there are canonical identifications
  \begin{align*}
    \TR^n_{\m\alpha,j}(S\rel\SAinf{R^+})
      &=
        \Fil^k \dRW[\ell]^j_{S/R^+}
        \tnsr
        \bigotimes_{i=0}^\ell \phi^{-i}(\xi)^{\m d_i}A\\
    \TR^n_{\m\alpha-1,j}(S\rel\SAinf{R^+})
      &=
        \frac{\dRW[k-1]^j_{S/R^+}}{\phi^{-k}(\xi)^{\m d_k}\dotsm\phi^{-\ell}(\xi)^{\m d_\ell}}
        \tnsr
        \bigotimes_{i=0}^{k-1} \phi^{-i}(\xi)^{\m d_i}A
  \end{align*}
  with $A=\Ainf(R^+$). This also holds for $n=\infty$ if $\ell<\infty$ or if $R^+=\O_C$.

  If $S$ is smooth over a perfect $\F_p$-algebra $k$, then there are canonical identifications
  \begin{align*}
    \TR^n_{\m\alpha,j}(S\rel\SW k)
      &= \Fil^{k-(\m d_k+\dotsb+\m d_\ell)} \dRW[\ell]^j_S\tnsr p^{\m d_0 + \dotsb + \m d_\ell}W\\
    \TR^n_{\m\alpha-1,j}(S\rel\SW k)
      &= \frac{\dRW[k-1]^j_S}{p^{\m d_k+\dotsb + \m d_\ell}}\otimes p^{\m d_0 + \dotsb + \m d_{(k-1)}}W
  \end{align*}
  with $W=\rmW(k)$. This also holds for $n=\infty$.
\end{theorem}

While the restriction to the contiguous range is disappointing, note that the contiguous range is already massively larger than the range of degrees that have been previously studied; besides, the point of this paper is to open the doors to more $\RO(G)$-graded ideas in the field, not to serve as the final word on it. Restricting to the contiguous range is most distressing for the following reason:
\[ \{\m\alpha \mid \m\alpha + i \text{ is contiguous for all } i \in \Z\} \]
is easily seen to be the union of the positive and negative ranges, so only in these ranges can we fully mix our result with the Postnikov filtration.

\begin{remark}
  The notations $\Z/n$ and $C_n$ are often used interchangeably for a cyclic group of order $n$. In this paper, we will think of $\Z/n$ as a quotient of $\widehat\Z$, while we think of $C_n$ as embedded in $\T$ (so perhaps a better notation would be $\Z/n\{-1\}$ as in Example \ref{ex:tor-ext}, the Pontryagin dual $(\Z/n)^\land$, or $\mu_n$, although the latter can be confused for a group scheme). As we will see in \sec\ref{sec:dRW}, the passage between these two perspectives involves the Cartier isomorphism (Example \ref{ex:cartier}); it would be interesting to know if there is a deeper connection between Pontryagin duality and the Cartier isomorphism.
\end{remark}

\begin{remark}[Poly-de Rham-Witt complex]
  Our computations imply a new expansion of the de Rham-Witt complex. In general, there is a map
  \[ \dRW[\bullet]^*_{S/R} \to \TR^\bullet_*(S \rel \SAinf R) \]
  because $\dRW[\bullet]^*$ is defined exactly as the universal example of ``the kind of structure $\TR^\bullet_*$ has''. This is quite vague, but in fact different authors get away with axiomatizing different fragments of the $d,F,V,R$ structure; and when $\TR$ becomes more complicated than the classical algebraic case, we can always generalize the definition of this structure, as with the $p$-typical Cartier complexes of \cite[Definition 3.34]{TCart} or the $\eta$-deformed ($p$-typical) Cartier complexes of \cite[Definition 4.17]{CycSyn}.

  The polyrepresentation-graded groups $\TR^\bullet_\rog$ possess $d$, $F$, $V$, and $R$ maps just as the $\Z$-graded groups do, but now also have $N$ maps. We can define the four variants ($p$-typical vs big, truncated vs full) of the \emph{poly-de Rham Witt complex}
  \begin{gather*}
    \dRW^\rog_{S/R}, \qquad \dRW[\bullet]^\rog_{S/R},\\
    \DRW^\rog_{S/R}, \qquad \DRW[\bullet]^\rog_{S/R}
  \end{gather*}
  as the universal examples of this structure. We will not study the poly-de Rham-Witt complex in this paper, but our computations provide an upper bound for it. Again, the poly-de Rham-Witt complex is a filtered and twisted version of the usual de Rham-Witt complex, rather than a totally new object, with the norms being the most interesting new piece of structure.
\end{remark}

\subsection{Norms}
While the $\RO(G)$-graded groups are interesting in their own right, our main reason for computing them is to be able to understand the Norm (even for nonzero integer degrees). Since the Norm is not very well-known, we begin by recalling the Norm map of Witt vectors from \cite{AngeltveitNorm,PrismsTambara1}.

\subsubsection{Norm map of Witt vectors}
This is a multiplicative function $N\colon \rmW(R) \to \rmW(R)$ characterized in ghost coordinates by
\begin{align*}
  FN(f) &= f^p\\
  N &\equiv \id \bmod V
\end{align*}
It is conceptually clarifying to also consider iterates of the Norm, characterized in ghost coordinates by
\begin{align*}
  F^k N^k(f) &= f^{p^k}\\
  F^{k-r} N^k(f) &\equiv f^{p^{k-r}} \bmod V, \quad 1\le r\le k\hspace{-5.1em}
\end{align*}
This descends to a map $N\colon \rmW_n(R) \to \rmW_{n+1}(R)$, in which case $N^n\colon R = \rmW_0(R) \to \rmW_n(R)$ is the familiar multiplicative lift, and in Witt coordinates
\begin{equation}
  \label{eq:witt-coords}
  (a_0, \dotsc, a_n) = N^n a_0 + VN^{n-1} a_1 + \dotsb + V^n a_n.
\end{equation}
This essentially says that $\rmW_\bullet(R)$ is the free $\Z_p$-Tambara functor on $R$, where $R$ is viewed as a ring with (trivial) $\Z_p$-action \cite{BrunWitt}.

Recall from \cite[Proposition 1]{JoyalDeltaWitt} and \cite[Lemma B.1]{DrinfeldFormalGroup} that a $\delta$-ring admits operations $\delta_\bullet$ and $\theta_\bullet$ satisfying
\begin{align*}
  \phi^n(f) &= \sum_{k=0}^n p^k \delta_k(f)^{p^{n-k}} &
  \phi(f^{p^{n-1}}) &= f^{p^n} + p^n \theta_n(f)
\end{align*}
(In particular, $\delta_0 = \id$ and $\delta_1 = \theta_1 = \delta$.) These give implicit and explicit formulas for iterates of the Norm:
\begin{align*}
  \id &= N + V\delta_1 &
  N &= \id - V\theta_1\\
  \id &= \sum_{k=0}^n V^k N^{n-k} \delta_k &
  N^k &= \id - \sum_{k=1}^n V^k \theta_k
\end{align*}
In particular, we can view $\id = \sum V^k N^{n-k} \delta_k$ as a ``Witt coordinates'' decomposition of the identity, generalizing \eqref{eq:witt-coords} to infinite level.

\subsubsection{de Rham-Witt norms}
Now let $S$ be $p$-completely smooth over a perfectoid ring $R$, let $A=\Ainf(R)$, and assume for the moment that $p>2$. By Theorem \ref{thm:intro-tr-ctg}, there are Norm maps that go
\begin{align}
  \label{eq:norm-dRW-map}
  N\colon \dRW[n]^j_{S/R} &\to \dRW[n+1]^{j}_{S/R} \fade{{}\tnsr_A ({}\dotsb)} + \Fil^1 \dRW[n+1]^{pj}_{S/R} \fade{{}\tnsr_A ({}\dotsb)}\\
  N^k\colon \dRW[n]^j_{S/R} &\to \sum_{r=0}^k \Fil^r \dRW[n+k]^{p^rj}_{S/R} \fade{{}\tnsr_A ({}\dotsb)}
\end{align}
where the sums are direct sums for $j>0$, and we are suppressing the ideal twists for simplicity; the precise twists are stated in Theorem \ref{thm:norms-dRW}. These are characterized in ghost coordinates by
\begin{align*}
  FN(\omega) &= \omega^p & F^k N^k(\omega) &= \omega^{p^k}\\
  N &\equiv \psi^p \bmod \Fil^1 & F^{k-r} N^k(\omega) &\equiv \psi^{p^r}(\omega)^{p^{k-r}} \bmod \Fil^1
\end{align*}

Here $\psi^\ell$ is the Adams operation $\psi^\ell\colon \Omega^*_{S/R} \to \Omega^*_{S/R}$ given by $\psi^\ell(\omega\in\Omega^j_{S/R}) = \ell^j \omega$. These are well-known for $p\nmid \ell$ (see \cite[Proposition 9.14]{BMS2} and \cite[Proposition 6.4.12]{RaksitDDR}), but we need the $p$-primary versions, which we study in \sec\ref{sub:adams}.

By reduction to the $p$-torsionfree case, it suffices to find a natural formula which satisfies the above congruences. We would like to try
\[ N(\omega) = \psi^p(\omega) - V\delta(\omega), \]
but it is not clear how to define $\delta$ on differential forms. Also, this does not have the right shape for the codomain of \eqref{eq:norm-dRW-map}.

The missing formula is
\begin{equation}
  \label{eq:norm-formula-2}
  N(\omega) =  (p-V(1))\frac{\psi^p(\omega)}p + V\frac{\omega^p}p
\end{equation}
It is easy to check that, as a formal expression, this satisfies the necessary congruences. This formula does not make sense in the degree 0 case, but it does in positive degree. Indeed, if $\omega$ is a differential form of positive degree, then $\omega^p$ and $\psi^p(\omega)$ are divisible by $p$, in fact canonically so in the $p$-torsionfree case, which we can reduce to by functoriality. Beware that the $V$ in $V\frac{\omega^p}p$ is an $\RO(G)$-graded Verschiebung.

With hindsight, we can manually extend $\psi^k$ to the de Rham-Witt complex via $\psi^k(\omega\in\dRW^j_{S/R}) = k^j\omega$, and then set
\[ \delta(\omega) \defeq \frac{F\psi^p(\omega) - \omega^p}p \]
which makes sense for the same reasons as above---though it is not a homogeneous operation! We can then rearrange \eqref{eq:norm-formula-2} into the unified formula
\[ \psi^p(\omega) = N(\omega) + V\delta(\omega) \]
which works also for $\omega\in\dRW^0$.

We can similarly extend $\delta_\bullet$ and $\theta_\bullet$ to the de Rham-Witt complex, and view the Norms along with Verschiebungs as providing a ``Witt coordinates'' decomposition of the Adams operations $\psi^{p^k}$.

\begin{warning}
Although this extension is conceptually appealing, the topological perspective only justifies the definition of $\psi^k$ up to $\Fil^1$ and only justifies $\delta$ up to $\ker V$.
\end{warning}

\begin{warning}
  It is not obvious that \eqref{eq:norm-formula-2} simplifies to $N(\omega)=V\frac{\omega^p}p$ in characteristic $p$, since some of the factor $(p-V(1))$ gets absorbed into the ideal twist. However, a small calculation shows that this factor still vanishes in characteristic $p$, at least for $p>2$ (Corollary \ref{cor:norms-Fp}).
\end{warning}

\begin{remark}
  \label{rmk:norms-2}
  When $p=2$, everything above goes through for de Rham-Witt forms of even degree. Producing Norms for odd-degree differential forms at $p=2$ requires contending with sign representations, which we defer to future joint work with Jason Yeung.
\end{remark}

\begin{remark}
  Let us rename the Norm discussed above (inherited from the Norm on $\TR$) to the \emph{topological Norm}, and denote it by $N_\top$. We define the \emph{algebraic Norm} $N_\alg$ by undoing the twists (see the explicit description of the twists below):
  \[\xymatrix{
      \dRW[n]^j_{S/R} \ar[r]^-{N_\top} \ar@{..>}[dr]_-{N_\alg} & \dRW[n+1]^j_{S/R} \tnsr ({}\cdots) + \Fil^1 \dRW[n+1]^{pj}_{S/R} \tnsr ({}\cdots) \ar[d]^-{\text{\tiny$\begin{bmatrix}1&0\\0&\xi^{j(p-1)/2}\end{bmatrix}$}}\\
                                       & \dRW[n+1]^j_{S/R} + \Fil^1 \dRW[n+1]^{pj}_{S/R} 
  }\]

  This is somewhat easier to work with, and we do not understand the significance of being able to refine $N_\alg$ to $N_\top$. For example, in characteristic $p$ it is obvious that $N_\alg(\omega) = V\frac{\omega^p}p$ for $j>0$; while this is not obvious for $N_\top(\omega)$, a nontrivial calculation shows that it is still true, so $N_\top$ is not seeing anything more. It is conceivable that $N_\top$ and $N_\alg$ will differ more significantly for odd-degree forms at $p=2$.
\end{remark}

\begin{example}
  \label{ex:norm-dx}
  Let $p$ be odd, and let $S$ be $p$-completely smooth over a perfectoid ring $R$.
  If $|\omega|=1$, then
  \begin{align*}
    N^k(\omega) &= (p^k - Vp^{k-1})\omega\\
    \delta_k(\omega) &= F^k\omega \fade{{} = \delta^k(\omega)}\\
    \theta_k(\omega) &= 0 \quad \text{ for } k>1
  \end{align*}
  Note that this implies the values of $N,\delta_k,\theta_k$ on all of $\dRW^*$ by the addition/multiplication formulas.

  Let us be more precise about the twist. Setting $h=\frac{p-1}2$, the Norm goes
  \[ N_{\fade\top} \colon \dRW[n]^1_{S/R} \to \dRW[n+1]^1_{S/R} \tnsr \phi^{-1}\left(\xi_{n+1}^h(h)_\prism!\right) \]
  The symbol $(i)_\prism!$ is defined in Definition \ref{dfn:prismatic-factorial}, and in the $q$-de Rham context the twist becomes
  \[ \phi^{-1}\left((p^{n+1})^h_{q^{1/p^{n+1}}} (h)_q!\right) = \phi^{-1}\prod_{i=1}^h (p^{n+1}j)_{q^{1/p^{n+1}}} \]
  In particular, the factor $(h)_{q^{1/p}}!$ means that the $F$ map out of the codomain is $h!$ times the usual map $F\colon \dRW[n+1]^1_{S/R} \to \dRW[n]^1_{S/R}$. But $FN\omega = \omega^p = 0$ anyway, so we do not understand the importance of this twist here.

  Ignoring the factorial unit, the refinement of $N_\alg$ to $N_\top$ amounts to saying that $p-V(1)$ is divisible by $\phi^{-1}\left(\xi_{n+1}^{(p-1)/2}\right)$ (Remark \ref{rmk:snell} and \cite[Lemma 2.1]{MaoLogarithm}). Taking $n=0$ for simplicity, this means that
  \[ N^1_{0,\top}(\omega) \approx \phi^{-1}(\xi)^{h+1}\omega \tnsr \phi^{-1}(\xi)^h \in \dRW[1]^1_{S/R} \tnsr \phi^{-1}\left(\xi^h (h)_\prism!\right) \]
  where $\approx$ means up to units.
  In characteristic $p>2$, we have $\frac{p+1}2 \ge \frac{3+1}2 = 2$, so the term $\phi^{-1}(\xi)^{h+1} = p^{\frac{p+1}2}$ vanishes in $\rmW_1(k)$, showing that indeed $N^1_{0,\top}(\omega)=0$.
\end{example}

\begin{example}
  \label{ex:norm-even}
  Let $S=R[x_1, y_1, \dotsc, x_p, y_p]^\land_p$ over an arbitrary perfectoid ring $R$. Let
  \[
    \tau = \sum_{i=1}^p d x_i\, d y_i \in \dRW^2_{S/R} \qquad
    \omega = \prod_{i=1}^p d x_i\, d y_i \in \dRW^{2p}_{S/R}
  \]
  We have $\tau^p = p!\omega$, so the algebraic Norm is
  \begin{align*}
    N_\alg(\tau)
      &= (p^2-V(p))\tau + V(p-1)!\omega\\
      &= V(p-1)!\omega \qquad\text{in characteristic $p$}
  \end{align*}
  Similarly,
  \begin{align*}
    \delta(\tau)
    &=
    \frac{F\psi^p(\tau) - \tau^p}p\\
    &= p\sum_{i=1}^p x_i^p y_i^p \dlog x_i \dlog y_i - (p-1)!\omega
  \end{align*}

  The topological Norm goes
  \[
    N_\top\colon
    \dRW[n]^2_{S/R}
    \to
    \bigoplus\left(
      \begin{array}{r@{}l}
      \dRW[n+1]^2_{S/R} &{}\tnsr \phi^{-1}\left(\xi_{n+1}^{p-1}(p-1)_\prism!\right)\\
      \Fil^1 \dRW[n+1]^{2p}_{S/R} &{}\tnsr \dfrac{\phi^{-1}(p-1)_\prism!}{\xi^{p-1}}
    \end{array}\right)
  \]
  so is given by
  \[
    N_\top(\tau) =
    \frac
      {p^2-V(p)}
      {\phi^{-1}\left(\xi_{n+1}^{p-1}(p-1)_\prism!\right)}
    \tau \fade{{}\tnsr \phi^{-1}\left(\xi_{n+1}^{p-1}(p-1)_\prism!\right)}
    +
    V\omega \fade{{}\tnsr\dfrac{\phi^{-1}(p-1)_\prism!}{\xi^{p-1}}} 
  \]
  (remembering that $F(1\tnsr\xi) = 1$, since $F$ is a \emph{divided} Frobenius.)
\end{example}

\subsubsection{Breuil-Kisin and prismatic norms}

Back to the Borel case, one would also like to compute the Tambara structure on $\THH(S\rel\S_A)^{\{h,t\}C_{p^\bullet}}_\rog$, extending the author's $\pi_0$ prismatic norms from \cite{PrismsTambara1}. Although we have managed to compute a few of these, they are far more complicated than expected, and we do not yet see an easy conceptual story like we have for the de Rham-Witt norms or $\pi_0$ prismatic norms. In the interest of this paper seeing the light of day, we postpone these to a future revision of this article. (If one restricts to perfectoid rings $R$, it is relatively easy to see that the norms on $\THH(R;\Z_p)^{\{h,t\}C_{p^\bullet}}_\rog$ take generators to generators, which suffices for many applications; the difficulty is in pinning down units.)

\subsection{Overview}

We review all our notation in \sec\ref{sec:notation}. We study the $\RO(G)$-graded homotopy of even Borel spectra in \sec\ref{sec:qrsp}: primarily the $\RO(C_{p^\infty})$-graded $\TCmin$ and $\TP$ of quasiregular semiperfectoid rings, and some parallel results for even $\E_\infty$-rings with $\T$-action. We study the de Rham-Witt setting in \sec\ref{sec:dRW}. We study $\RO(G)$-graded norms in \sec\ref{sec:norms}.

A web interactive for exploring the homotopy orbits to TR spectral sequence is available at

\begin{center}
  \url{https://ysulyma.github.io/papers/poly/}
\end{center}

\subsection{Acknowledgments}

We thank Benjamin Antieau, Andrew Blumberg, Elden Elmanto, Lars Hesselholt, Jacob Lurie, Zhouhang Mao, Shubhodip Mondal, and Noah Riggenbach for helpful conversations related to this work. We are especially grateful to Antieau for discussions about the Adams operation $\psi^p$; Mao for explaining Example \ref{ex:tor-ext} to us; and to Riggenbach for clarifying the difference between absolute and relative Ngyaard filtrations for us.
